\documentclass{article}

\usepackage[english]{babel}
\usepackage[letterpaper,top=2cm,bottom=2cm,left=3cm,right=3cm,marginparwidth=1.75cm]{geometry}

\usepackage{amsmath}
\usepackage{graphicx}
\usepackage{xcolor}
\usepackage[colorlinks=true, allcolors=blue]{hyperref}
\usepackage{subcaption}
\usepackage{tabularray}
\usepackage{float}
\usepackage[normalem]{ulem}

\UseTblrLibrary{booktabs}
\UseTblrLibrary{siunitx}

\newcommand{\tinytableTabularrayUnderline}[1]{\underline{#1}}
\newcommand{\tinytableTabularrayStrikeout}[1]{\sout{#1}}
\NewTableCommand{\tinytableDefineColor}[3]{\definecolor{#1}{#2}{#3}}

\title{Problem Set 9}
\date{April 14, 2026}
\author{Caleb Reid}

\begin{document}

\maketitle

The optimal value of $\lambda$  for the LASSO model, identified by 6-fold cross validation, was 0.00356. The in-sample and out-sample Root Mean Square Errors (RMSE) were 0.0735 and 0.172, respectively.

For the ridge model, the penalty parameter or $\lambda$ was equal to 0.0233. In-sample RMSE was equal to 0.0687 and the out-of-sample RMSE was 0.173.

It is not possible to estimate a simple linear regression model when the number of columns is greater than the number of observations or K \textgreater  n. In terms of the bias-variance tradeoff, both models are in the "Goldilocks zone" of the bias-variance curve. The models performed similarly on the test or out-of-sample data as they did on the training or in-sample data.  The low RMSE values show that the models are not overfitting on the in-sample data.

\end{document}